\chapter{La variable objetivo}
\label{cap:variable}

% BLOQUE «Planteamiento de la comparativa» (3/3). Es LA MEDIDA sobre la que se compara:
% la norma del tipo 3 exige identificarla en este bloque, antes del desarrollo.

% GUION. CAPITULO DIFERENCIAL: es el resultado central del trabajo.
% Fuente: TFM-Quantum/doc/hallazgo_objetivo.md + doc/memoria/materiales.md seccion 4.
% APROBADO POR EL DIRECTOR (ago-2026).
%
% ATENCION: todas las cifras de este capitulo son del dataset v1. HAY QUE REHACERLAS
% con el v2 antes de escribirlas. Ver doc/memoria/pendientes.md.


\section{Los observables bajo ruido}
\label{sec:observables_ruido}
% GUION. Seccion de ENTRADA del capitulo: antes de discutir QUE se predice, hay que
% enseñar QUE le pasa a un observable cuando el circuito se ejecuta con ruido.
% Sin esto, el lector llega a la definicion de Delta sin saber de que magnitud se habla.
%
% Fuentes: doc/02_Los_Observables.md · doc/memoria/materiales.md §4 y §5 ·
%          TFM-Quantum/doc/changes/2026-08_cambio_decisiones.md
%
% 🔴 TODAS LAS CIFRAS DE ABAJO SON DEL v1 O ESTAN PENDIENTES DE REHACER SOBRE EL v2.
%    NO escribir ninguna sin recalcularla. Ver doc/memoria/pendientes.md §1.

\subsection{Qué le hace el ruido a un valor esperado}
% La intuicion de partida, y por que es incompleta:
%   - la idea ingenua es que el ruido ATENUA: acerca el valor esperado a cero, porque
%     destruye la coherencia que lo sostenia;
%   - por eso la magnitud natural para describirlo es un COCIENTE y no una diferencia:
%     el FACTOR DE SUPERVIVENCIA f = <O>ruidoso / <O>exacto, la fraccion de señal que
%     sobrevive a la ejecucion;
%   - f = 1 significa que el circuito salio intacto; f = 0, que no queda nada.
% ⚠️ Enlazar con el apartado~\ref{sec:puertas}: cuantas mas puertas, menor f.

\subsection{Reducción y amplificación}
% EL PUNTO DE LA SECCION, y lo que la hace necesaria: el ruido NO SIEMPRE ENCOGE.
% Que contar:
%   - el caso mayoritario es la REDUCCION (f < 1): la señal se atenua hacia cero;
%   - pero hay AMPLIFICACION (f > 1) en una fraccion no despreciable de las muestras,
%     y hasta CAMBIO DE SIGNO (f < 0);
%   - y donde ocurre no es aleatorio: pasa casi siempre cuando el valor exacto es
%     DIMINUTO, porque entonces el cociente se dispara aunque la diferencia absoluta
%     sea minuscula. Es un artefacto de dividir por algo cercano a cero, no un
%     fenomeno fisico nuevo.
%
% 🔴 CIFRAS PENDIENTES DE REHACER SOBRE EL v2 (las del v1 no son citables):
%     - que porcentaje de muestras tiene f > 1  (v1: «falla hasta el 15 % de las veces»);
%     - el rango real de f y la cola: tras retirar las tres dispersiones, f pasa de
%       alcanzar valores de 71 a quedarse en [-0,9 · 2,3], y las filas con |f| > 2 caen
%       del ~20 % a menos del 1 %;
%     - que fraccion de muestras NO TIENE SEÑAL que perder (v1: entre el 22 % y el 37 %).
%
% ⚠️ Una figura aqui valdria mucho: histograma de f, o dispersion de <O>ruidoso frente a
%    <O>exacto con la diagonal f=1 marcada. Se ve de un vistazo que la nube cae por
%    debajo de la diagonal pero que hay puntos por encima.

\subsection{El ruido no solo atenúa: también desplaza}
% Matiz MEDIDO que corrige la intuicion del apartado anterior. El modelo real no es
% multiplicativo puro sino AFIN:
\begin{equation}
\label{eq:modelo_afin}
\langle O \rangle_{\text{ruidoso}} = f \cdot \langle O \rangle_{\text{exacto}} + b .
\end{equation}
% Que contar:
%   - `b` es un desplazamiento que no depende de cuanta señal habia;
%   - medido: en los cinco observables vigentes aporta entre el 1 % y el 17 % de la
%     correccion, y es NEGATIVO en tres de los cinco;
%   - DECISION: `b` NO se modela. El modelo predice solo `f`. Justificarlo con esa cifra.
%   - y este es ademas el criterio que dejo fuera a las tres dispersiones: alli `b` valia
%     el 135-149 % de la propia degradacion (capitulo~\ref{cap:dataset}).

\subsection{Un aviso sobre el rango de \texorpdfstring{$f$}{f}}
% 🔴 NUNCA ESCRIBIR que f esta entre 0 y 1 «por fisica». Es falso y comprobado:
%    el 6,4 % de las etiquetas reales cae fuera de [0,1].
%    La prediccion se recorta a [0,05 · 1,0] por ESTABILIDAD NUMERICA —para que el
%    cociente de la reconstruccion no explote—, no porque la fisica lo impida.
%    Confundir las dos cosas es justo el tipo de afirmacion que un tribunal comprueba.

\section{Lo que se quería predecir}
\label{sec:delta}
El objetivo natural: $\Delta = \langle O \rangle_{\text{exacto}} - \langle O
\rangle_{\text{ruidoso}}$, la degradación que sufre el observable. Predecirla permitiría
restarla del resultado medido.

\section{La medida: no es predecible}
\label{sec:no_predecible}
% Tabla de R2 sobre los tres objetivos. REHACER con el v2.
Se ajustó \textbf{el mismo modelo, con las mismas variables y la misma partición}, cambiando
únicamente el objetivo.

\textbf{$\Delta$ con signo da $R^2 \approx 0$ en validación.} Y el detalle importa: en
\textbf{validación}, no en test. No es un problema de generalización, ni de arquitectura, ni
de cantidad de datos.

\section{Por qué: la causa es aritmética}
\label{sec:causa}
\begin{equation}
\label{eq:delta_factoriza}
\Delta = (1 - f)\cdot\langle O \rangle_{\text{exacto}}, \qquad
f = \frac{\langle O \rangle_{\text{ruidoso}}}{\langle O \rangle_{\text{exacto}}}
\end{equation}

$\Delta$ es el producto de dos cosas de naturaleza distinta:
\begin{itemize}
    \item $(1-f)$ ( \textbf{cuánta señal destruye el ruido}. Es física del hardware: depende
          del circuito y de la calibración, y por tanto \textbf{es aprendible}.
    \item $\langle O \rangle_{\text{exacto}}$ ) \textbf{cuánta señal había}. Depende del
          algoritmo, y es justamente la incógnita para la que haría falta un ordenador
          cuántico. \textbf{No es aprendible.}
\end{itemize}

\subsection{La excepción que confirma la explicación}
% Argumento fuerte: no depender solo de la formula.
Hay un observable cuyo $\Delta$ tiene signo sistemático (positivo el $99{,}9$\,\% de las veces).
\textbf{Es el único con $R^2$ positivo sobre $\Delta$}: donde no hay signo que adivinar, el
problema se vuelve aprendible.

\section{La reformulación}
\label{sec:reformulacion}
Predecir $f$ y reconstruir $\Delta$ con el valor \textbf{medido}:
\begin{equation}
\label{eq:reconstruccion}
\hat\Delta = \langle O \rangle_{\text{ruidoso}}\cdot\frac{1 - \hat f}{\hat f}
\end{equation}

\subsection{Por qué NO rompe el paradigma pre-ejecución}
% ESTE es el punto que hay que defender ante el tribunal. Redactarlo con cuidado.
\begin{itemize}
    \item \textbf{La entrada del modelo no cambia}: sigue siendo solo el circuito y la
          telemetría del chip. El modelo nunca ve el resultado de la ejecución.
    \item \textbf{El valor medido entra solo en la aritmética final} de la corrección —
          exactamente donde ya entraba en $\langle O \rangle_{\text{mitigado}} = \langle O
          \rangle_{\text{ruidoso}} + \Delta$.
          % ⚠️ EL SIGNO ES UNA SUMA, y es un error facil de arrastrar. Delta esta definido
          %    en el dataset como Delta = exacto − ruidoso, luego exacto = ruidoso + Delta.
          %    Con una resta, un modelo perfecto DUPLICARIA el error en vez de anularlo.
          %    Corregido en sep-2026; hasta entonces cinco documentos ponian la resta.
          %    Es coherente con la ecuacion~\ref{eq:reconstruccion}: si ruidoso = f·exacto,
          %    entonces Delta = ruidoso·(1−f)/f, que es POSITIVO cuando f < 1.
    \item QEMFormer necesita el valor ruidoso \textbf{como variable de entrada del modelo}.
          La diferencia sigue siendo real.
\end{itemize}

\subsection{Una ventaja añadida}
El caso degenerado se resuelve solo: si el valor medido es $\approx 0$, entonces
$\hat\Delta \approx 0$ — que es la respuesta correcta cuando no hay señal que corregir. En la
formulación original ese caso era ruido puro, y afecta a una fracción grande del conjunto.

\section{El límite honesto de la solución}
\label{sec:limite}
% NO adornar esto. Es lo que hace creible el resto.
% ⚠️ CIFRAS DEL v1 y sobre OCHO observables. El objetivo vigente son CINCO
%    (se retiraron las tres dispersiones, ago-2026). Rehacer antes de escribirlo.
En validación el esquema funciona en todos los observables. \textbf{En test solo transfiere
uno.} En el resto, el $R^2$ se hunde pese al buen ajuste en validación.

\textbf{Las variables agregadas aprenden la física, pero no la transfieren a tipos de circuito
nunca vistos.}

\subsection{Por qué esto no es un fracaso}
Es el resultado que el conjunto de datos fue \textbf{diseñado para medir}, y deja a la
comparativa del capítulo~\ref{cap:comparativa} una pregunta real:

\begin{quote}
¿Transfiere la \textbf{red de grafos}, que ve la estructura del circuito, donde el
resumen agregado no lo hace?
\end{quote}

Con $\Delta$ como objetivo único, los tres modelos habrían dado $R^2 \approx 0$ en validación
\textbf{y} en test, y la comparativa habría sido entre tres formas de no predecir nada.

\section{Alternativas descartadas}
\label{sec:alternativas_objetivo}
Mantener $\Delta$ (sin señal que comparar), eliminar $\Delta$ de la memoria (pierde la métrica
física que conecta con la literatura), o añadir el valor ruidoso como entrada (es el enfoque
de QEMFormer y elimina el diferenciador del trabajo).
